\documentclass[9pt]{beamer}
\usetheme{metropolis}
\usepackage{amsmath}
\usepackage{amssymb}
\usepackage{graphicx}
\usepackage{booktabs}
\usepackage{xcolor}

% Define color palette matching our project design system
\definecolor{primary}{HTML}{1abc9c}    % Deep Teal
\definecolor{secondary}{HTML}{3498db}  % Astro Blue
\definecolor{darkbg}{HTML}{2c3e50}     % Slate Blue
\definecolor{accent}{HTML}{e67e22}     % Orange

\setbeamercolor{palette primary}{bg=darkbg, fg=white}
\setbeamercolor{palette secondary}{bg=primary, fg=white}
\setbeamercolor{palette tertiary}{bg=secondary, fg=white}
\setbeamercolor{structure}{fg=primary}
\setbeamercolor{alerted text}{fg=accent}

\title{Stellar Trace: Physics-Informed Modeling \& Bayesian Inference}
\subtitle{Meet 3: Bayesian MCMC, PINN Convolutions, \& Delay Time Distributions (Tasks 13--15)}
\author{Astronomy Club}
\date{\today}

\begin{document}

\maketitle

% Slide 2: Meet Series Roadmap
\begin{frame}{Meet Series Roadmap}
  \begin{block}{Phase 2 Curriculum Structure}
    Stellar Trace Phase 2 transitions from the volume-complete universe model of Milestone 1 to real data ingestion, classifier training, and physical model fitting:
  \end{block}
  \begin{description}
    \item[Meet 1:] \textbf{Telescopes \& Statistics (Tasks 7--9)}
      \begin{itemize}
        \item Selection effects, coordinate cross-matching, database querying, and 1D/2D K-S tests.
      \end{itemize}
    \item[Meet 2:] \textbf{Machine Learning \& Data Leakage (Tasks 10--12)}
      \begin{itemize}
        \item Stratified splits, manual random oversampling, MLP classifiers, GridSearchCV, and Permutation Feature Importance.
      \end{itemize}
    \item[Meet 3 (Today):] \textbf{Bayesian MCMC \& PINN Convolutions (Tasks 13--15)}
      \begin{itemize}
        \item Task 13: Physics-Informed Supernova Rate Modeling.
        \item Task 14: Model Interpretability and Physical Comparison.
        \item Task 15: Type Ia Supernova Delay Time Distribution (DTD) Reconstruction.
      \end{itemize}
  \end{description}
\end{frame}

% Slide 3: Meet 3 Objectives (Part 1)
\begin{frame}{Meet 3 Objectives (Part 1)}
  By the end of this meet, you will be able to:
  \begin{enumerate}
    \item Define competing astrophysical models for supernova rates as functions of host galaxy properties.
    \item Derive the likelihood and Negative Log-Likelihood (NLL) functions for host galaxy observations.
    \item Solve for point-estimate physical parameters using classical Maximum Likelihood Estimation (MLE).
    \item Implement a Metropolis-Hastings Markov Chain Monte Carlo (MCMC) sampler to quantify parameter uncertainties and map degeneracies.
    \item Construct a Physics-Informed Neural Network (PINN) to recover physical parameters via gradient descent.
  \end{enumerate}
\end{frame}

% Slide 4: Meet 3 Objectives (Part 2)
\begin{frame}{Meet 3 Objectives (Part 2)}
  By the end of this meet, you will also be able to:
  \begin{enumerate}
    \setcounter{enumi}{5}
    \item Validate rate models using 2-sample Kolmogorov-Smirnov (K-S) goodness-of-fit validation tests.
    \item Compare the interpretability of physically grounded parameters with black-box feature importances.
    \item Formulate and discretize integral convolutions of Star Formation History (SFH) with the Delay Time Distribution (DTD).
    \item Solve convolved DTD systems in PyTorch to reconstruct the power-law slope $\gamma$ of white dwarf mergers.
  \end{enumerate}
\end{frame}

% Slide 4: Section 1 Divider
\begin{frame}[plain]
  \vfill
  \centering
  \begin{beamercolorbox}[sep=8pt,center,shadow=true,rounded=true]{palette primary}
    \usebeamerfont{title}\textbf{Section 1: Physics-Informed Supernova Rate Modeling (Task 13)}
  \end{beamercolorbox}
  \vfill
\end{frame}

% Slide 5: Competing Astrophysical Rate Models
\begin{frame}{Competing Astrophysical Rate Models}
  In previous tasks, we compared host galaxy properties with background populations using statistical tests. Now, we define three competing models for supernova rates as functions of host galaxy properties:
  \begin{description}
    \item[1. Null (Random) Model:] Supernovae are random events that do not depend on host properties:
      $$R_{\text{Null}} = \text{const}$$
    \item[2. Core-Collapse Supernova (CC-SN) Model:] Tracks instantaneous star formation as a power law of the Star Formation Rate (SFR):
      $$R_{\text{CC}}(\text{SFR}) = \text{SFR}^\beta$$
      where $\beta \approx 1.0$ if the rate is directly proportional to the active star formation.
    \item[3. Type Ia Supernova Model:] The standard ``A+B'' model with delayed (stellar mass) and prompt (recent star formation) components:
      $$R_{\text{Ia}}(M_\star, \text{SFR}) = A \cdot \left(\frac{M_\star}{10^{10} M_\odot}\right) + B \cdot \left(\frac{\text{SFR}}{M_\odot \text{ yr}^{-1}}\right)$$
      To simplify calculations and ensure numerical stability, we define:
      $$R_{\text{Ia}}(M_\star, \text{SFR}) = A \cdot M_{\star, 10} + B \cdot \text{SFR}$$
  \end{description}
\end{frame}

% Slide 6: Parameter Recovery via Maximum Likelihood Estimation (MLE)
\begin{frame}{Parameter Recovery via Maximum Likelihood Estimation (MLE)}
  The probability of a specific galaxy $i$ hosting a supernova out of a population of $N$ galaxies is proportional to its relative rate:
  $$P(\text{host}_i \mid \theta) = \frac{R_i(\theta)}{\sum_{j=1}^N R_j(\theta)}$$
  For a sample of observed host galaxies, the joint likelihood is the product of their individual probabilities.
  \begin{block}{Negative Log-Likelihood (NLL) Formulation}
    To find the optimal parameters $\theta$, we minimize the Negative Log-Likelihood ($\mathcal{L}$):
    \begin{align*}
      \mathcal{L}(\theta) &= - \ln \prod_{i \in \text{observed}} P(\text{host}_i \mid \theta) \\
      &= - \sum_{i \in \text{observed}} \ln \left( \frac{R_i(\theta)}{\sum_{j \in \text{background}} R_j(\theta)} \right) \\
      &= -\sum_{i \in \text{observed}} \ln R_i(\theta) + N_{\text{observed}} \cdot \ln \left(\sum_{j \in \text{background}} R_j(\theta)\right)
    \end{align*}
  \end{block}
\end{frame}

% Slide 7: Classical MLE Fit Workflow
\begin{frame}{Classical MLE Fit Workflow}
  To fit our physical parameters ($\beta$ for CC-SNe; $A, B$ for Type Ia SNe) using point-estimate optimization:
  \begin{itemize}
    \item \textbf{Data Extraction:} Compile the observed host parameters ($M_\star, \text{SFR}$) and the background field galaxy parameters from the detected catalog.
    \item \textbf{L-BFGS-B Optimization:} Pass the custom NLL functions to a bounded minimizer (e.g., \texttt{scipy.optimize.minimize}).
    \item \textbf{Boundary Constraints:}
      \begin{itemize}
        \item Enforce parameter bounds: $\beta \in [0.1, 3.0]$ to prevent physical collapse.
        \item Enforce non-negativity: $A, B \in [0.0, 5.0]$ since supernova rates cannot be negative.
      \end{itemize}
    \item \textbf{Output Estimates:} Recover point estimates, showing that CC-SNe trace SFR ($\beta \approx 1.0$), and Type Ia SNe depend on both mass and star formation ($A, B > 0$).
  \end{itemize}
\end{frame}

% Slide 8: Bayesian Uncertainty Quantification via MCMC
\begin{frame}{Bayesian Uncertainty Quantification via MCMC}
  Maximum Likelihood Estimation gives only point estimates. In astrophysics, parameter degeneracies are common (e.g., mass and SFR are correlated, causing degeneracies between $A$ and $B$).
  To map the full posterior probability distribution and measure uncertainties, we implement a \textbf{Metropolis-Hastings MCMC sampler}:
  \begin{block}{Metropolis-Hastings Algorithm Step-by-Step}
    \begin{enumerate}
      \item **Initialization:** Start at initial parameter state $\theta_t = [A_0, B_0]$.
      \item **Proposal:** Draw a candidate step from a symmetric Gaussian distribution:
        $$\theta_{\text{prop}} \sim \mathcal{N}(\theta_t, \sigma^2)$$
      \item **Prior Check:** Enforce flat prior boundaries $A, B \in [0, 5]$. If candidate falls outside, reject and repeat current state.
      \item **Acceptance:** Calculate the acceptance ratio:
        $$\alpha = \min\left(1, \exp(\ln \mathcal{L}(\theta_{\text{prop}}) - \ln \mathcal{L}(\theta_t))\right)$$
      \item **Update:** Draw $u \sim \mathcal{U}(0, 1)$. If $u < \alpha$, accept $\theta_{t+1} = \theta_{\text{prop}}$; else, reject $\theta_{t+1} = \theta_t$.
    \end{enumerate}
  \end{block}
\end{frame}

% Slide 9: Physics-Informed Neural Network (PINN) Parameter Recovery
\begin{frame}{Physics-Informed Neural Network (PINN) Parameter Recovery}
  Instead of classical numerical solvers, we can recover parameters using a PyTorch neural network that incorporates our physical rate equations:
  \begin{itemize}
    \item \textbf{Constraining Parameters:} We define parameters ($\beta, A, B$) as neural network weights. To enforce strict positive boundaries ($A, B, \beta > 0$) without coordinate clipping, we parameterize their raw inputs and apply an exponential layer:
      $$\beta = e^{\tilde{\beta}}, \quad A = e^{\tilde{A}}, \quad B = e^{\tilde{B}}$$
    \item \textbf{Forward Propagation:} The network takes the host properties as inputs and computes the rates directly using the physical equations:
      $$R_{\text{Ia}} = A \cdot M_{\star, 10} + B \cdot \text{SFR}$$
    \item \textbf{Optimization Loss:} The model minimizes the joint Negative Log-Likelihood loss using the Adam optimizer. Gradients flow from the loss directly back to the physical weights ($\tilde{\beta}, \tilde{A}, \tilde{B}$) via automatic differentiation.
  \end{itemize}
\end{frame}

% Slide 10: Model Validation \& Goodness-of-Fit
\begin{frame}{Model Validation \& Goodness-of-Fit}
  How do we prove that our physical models are superior to the null model?
  \begin{itemize}
    \item \textbf{Synthetic Catalog Generation:}
      Using the fitted MLE parameters, we calculate rate weights $w_i = R_i / \sum R_j$ for every galaxy in the detected mock universe catalog. We sample 5,000 galaxies with replacement using these weights.
    \item \textbf{Goodness-of-Fit Testing:}
      We run 2-sample Kolmogorov-Smirnov (K-S) tests on the physical distributions of the simulated catalogs against the observed hosts:
      \begin{itemize}
        \item **Null Model Comparison:** Compares observed hosts against randomly sampled non-host galaxies. Yields $p$-values $\ll 10^{-5}$, rejecting the null hypothesis.
        \item **Physics Model Comparison:** Compares observed hosts against physics-weighted mock galaxies. Yields $p$-values $> 0.05$ (e.g., $p \approx 0.50$), confirming we cannot distinguish our model's predictions from reality.
      \end{itemize}
  \end{itemize}
\end{frame}

% Slide 11: Section 2 Divider
\begin{frame}[plain]
  \vfill
  \centering
  \begin{beamercolorbox}[sep=8pt,center,shadow=true,rounded=true]{palette primary}
    \usebeamerfont{title}\textbf{Section 2: Model Interpretability \& Comparison (Task 14)}
  \end{beamercolorbox}
  \vfill
\end{frame}

% Slide 12: Black-Box MLP vs. Physical PIML
\begin{frame}{Black-Box MLP vs. Physical PIML}
  Let's contrast the scientific value of our two Phase 2 approaches:
  \begin{columns}
    \begin{column}{0.5\textwidth}
      \alert{\textbf{MLP Classifier (Tasks 11--12):}}
      \begin{itemize}
        \item **Output:** Class probabilities mapping galaxy inputs to classifications.
        \item **Interpretability:** Permutation Feature Importance shows that mass and SFR are highly important features.
        \item **Limitations:** Does not yield structural physical laws or rate coefficients. Output is a black-box correlation.
      \end{itemize}
    \end{column}
    \begin{column}{0.5\textwidth}
      \alert{\textbf{PIML Rate Model (Task 13):}}
      \begin{itemize}
        \item **Output:** Physically grounded rate functions $R(M_\star, \text{SFR})$.
        \item **Interpretability:** Solve for physical constants: $\beta \approx 1.0$ (proves CC-SNe trace massive star formation); $A, B > 0$ (quantifies prompt vs. delayed Type Ia supernova rates).
        \item **Advantages:** Yields parameter uncertainties, maps covariance degeneracies, and enforces physical laws.
      \end{itemize}
    \end{column}
  \end{columns}
\end{frame}

% Slide 13: Section 3 Divider
\begin{frame}[plain]
  \vfill
  \centering
  \begin{beamercolorbox}[sep=8pt,center,shadow=true,rounded=true]{palette primary}
    \usebeamerfont{title}\textbf{Section 3: Delay Time Distribution (DTD) Reconstruction (Task 15)}
  \end{beamercolorbox}
  \vfill
\end{frame}

% Slide 14: Progenitor Delay Times & Delay Time Distributions
\begin{frame}{Progenitor Delay Times \& Delay Time Distributions}
  Type Ia supernovae originate from binary systems where a carbon-oxygen white dwarf reaches the Chandrasekhar mass limit or undergoes a double white dwarf merger.
  \begin{itemize}
    \item \textbf{Delay Time ($\tau$):} The time elapsed between the formation of the star system (starburst) and the subsequent supernova explosion.
    \item \textbf{Delay Time Distribution ($\Psi(\tau)$):} The distribution of delay times for a single-burst stellar population.
    \item \textbf{Progenitor Power Law:} Theoretical white-dwarf merger models suggest that the DTD follows a power law:
      $$\Psi(\tau) \propto \tau^{-\gamma}$$
      where $\gamma \approx 1.0 - 1.2$ for delay times $\tau \ge t_{\text{min}} \approx 100\text{ Myr}$.
  \end{itemize}
\end{frame}

% Slide 15: Convolution with Star Formation History (SFH)
\begin{frame}{Convolution with Star Formation History (SFH)}
  A real galaxy is not a single stellar population. It forms stars continuously over cosmic time.
  \begin{block}{Convolved Rate Model}
    The supernova rate in a galaxy at age $t$ is the convolution of its historical Star Formation History ($\text{SFR}(t)$) with the DTD ($\Psi$):
    $$R_{\text{Ia}}(t) = \int_0^t \text{SFR}(t - \tau) \Psi(\tau) d\tau = \int_0^t \text{SFR}(t - \tau) \tau^{-\gamma} d\tau$$
  \end{block}
  \begin{block}{Analytical Star Formation History Model}
    We approximate the SFH using a standard delayed-exponential model:
    $$\text{SFR}(t \mid \tau_{\text{SF}}) = \text{SFR}_0 \cdot \frac{t}{\tau_{\text{SF}}} e^{-t/\tau_{\text{SF}}}$$
    where $\tau_{\text{SF}}$ is the characteristic star formation timescale.
  \end{block}
\end{frame}

% Slide 16: Discretizing the Convolution Integral
\begin{frame}{Discretizing the Convolution Integral}
  To evaluate convolved rates numerically, we map each galaxy to a cosmic time grid $t_k = k \cdot \Delta t$ from $t=0$ to the galaxy's age $T \approx 13.7\text{ Gyr}$:
  \begin{block}{Discretized Rate Equation}
    $$R_{\text{Ia}}(T \mid \tau_{\text{SF}}, \gamma) = \sum_{\tau_k \ge t_{\text{min}}}^{T} \text{SFR}(T - \tau_k \mid \tau_{\text{SF}}) \cdot \tau_k^{-\gamma} \cdot \Delta t$$
  \end{block}
  \begin{itemize}
    \item We define the delay time grid step $\Delta t = 0.1\text{ Gyr}$.
    \item $\text{SFR}(T - \tau_k \mid \tau_{\text{SF}})$ represents the star formation rate of the galaxy at the lookback time $\tau_k$.
    \item $\tau_k^{-\gamma}$ represents the fraction of stars born at lookback time $\tau_k$ that explode today.
    \item The total rate is the sum over all lookback time steps.
  \end{itemize}
\end{frame}

% Slide 17: Reconstructing $\gamma$ via PyTorch Optimization: The Solver
\begin{frame}{Reconstructing $\gamma$ via PyTorch: The Solver}
  We recover the power-law slope $\gamma$ from observed hosts by optimizing a PyTorch convolved rate solver:
  \begin{itemize}
    \item \textbf{Parameterization:} Define $\gamma$ as an exponential parameter to enforce physical boundaries: $\gamma = e^{\tilde{\gamma}}$.
    \item \textbf{PyTorch Solver:} Evaluates the vectorized numerical convolution across the training galaxies:
      \begin{enumerate}
        \item Computes the $\text{SFR}$ grid matrix for all galaxies based on their age and $\tau_{\text{SF}}$.
        \item Computes the DTD vector $\boldsymbol{\psi}(\gamma) = \tau^{-\gamma}$.
        \item Performs matrix multiplication to evaluate rates: $\mathbf{r} = \mathbf{S} \boldsymbol{\psi} \cdot \Delta t$.
      \end{enumerate}
  \end{itemize}
\end{frame}

% Slide 18: Reconstructing $\gamma$ via PyTorch Optimization: The Loss
\begin{frame}{Reconstructing $\gamma$ via PyTorch: The Loss}
  To train the convolved solver, we optimize the network weights via gradient descent:
  \begin{itemize}
    \item \textbf{Loss Function Optimization:} Minimizes the Negative Log-Likelihood:
      $$\mathcal{L} = -\sum_{i \in \text{hosts}} \ln R_i + N_{\text{hosts}} \ln \left( \sum_{j \in \text{all}} R_j \right)$$
      using the Adam optimizer.
    \item \textbf{Backpropagation:} Gradients flow from the loss function through the numerical convolution layer back to $\tilde{\gamma}$ via PyTorch's autograd.
    \item \textbf{Result:} The convolved solver successfully recovers the true power-law slope $\gamma \approx 1.15$.
  \end{itemize}
\end{frame}

% Slide 18: Summary of Phase 2
\begin{frame}{Summary of Phase 2}
  \begin{block}{Key Milestones Completed}
    \begin{itemize}
      \item **Meet 1 (Tasks 7--9):** Applied observational selection filters (apparent magnitude cuts) to model Malmquist Bias, queried SDSS, and performed 1D/2D K-S tests.
      \item **Meet 2 (Tasks 10--12):** Resolved class imbalance without data leakage, trained MLP classifiers, and analyzed feature importances.
      \item **Meet 3 (Tasks 13--15):** Modeled physical supernova rate equations, mapped parameter uncertainties using MCMC, and reconstructed Delay Time Distributions via PyTorch convolved rate solvers.
    \end{itemize}
  \end{block}
  \begin{exampleblock}{Astrophysical Conclusion}
    We have successfully integrated machine learning with astrophysics: moving from statistical tests to predicting host suitability, and finally extracting physical laws using Physics-Informed models.
  \end{exampleblock}
\end{frame}

\end{document}
